% BEGIN THERMAL_R03_SYNC_13
\section{Conclusions}
\label{sec:ii-conclusiones}

La construction obtient les sections d’action de Planck et leur quotient
de retour, évalue le paramètre de Barbero--Immirzi à partir de l’incidence
orientée et produit les angles CKM avec leur phase et leurs amplitudes unitaires.
Elle détermine également le transducteur de Boltzmann pour une section thermique
positive et le coefficient gravitationnel commun de la réalisation R1--R8,
dont la spécialisation circulaire conserve la coïncidence des rayons.
L’inscription détermine \(r_N=5759/23040\) ; la composition longitudinale
donne \(L=r_N\alpha_{\HMT}^{16}L_*\) ; la métrique radiale et l’action
fixent \(G_{\rm ret}=c_{\rm int}^3L^2/\hbar_{\rm ret}\) sous
l’identification explicite du rayon réduit. Les identités résultantes relient les constantes
à leurs fonctions : énergie par fréquence, résolution élémentaire d’aire,
information de mélange, conversion thermique et invariance géométrique.

La relation centrale consiste dans la conservation de l’échelle géométrique.
Le produit \(G_a\hbar_a\) conserve l’échelle aréale tandis que l’action
et la gravitation se transforment réciproquement ; l’incidence pondérée
par \(\gamma_{\HMT}\) détermine la résolution de cette aire ; et
\(k_B\Theta_{{\rm clk},a}\log3=h_a/T_{108}\) relie la même action
à l’énergie d’une décision tritique. L’observation conjointe de l’aire
et du hamiltonien modulaire, avec \(\Delta_{\rm mod}\ne0\) et une
normalisation connue, retrouve les occupations sectorielles qu’une
seule lecture agrégée ne distingue pas. L’inversion directe de
\((N,D)\) reste valable sans cette condition. Chaque identité conserve les
conditions de sa réalisation dimensionnelle, géométrique ou spectrale.

\subsection{Constantes de Planck, Barbero--Immirzi et mélange CKM}

L’échelle d’action contient deux déterminations distinctes. Le
quotient local du registre a seize feuillets ; la norme multiplicative
et l’auto-échelle fixent \(\operatorname{ord}_{10}(\varphi\alpha^{16})=-34\). Les caractères régionaux fixent les sections orientées
et le quotient de retour. La conversion entre action angulaire et action
par tour maintient l’identité \(h=2\pi\hbar\). La coordinatisation angulaire
transporte les représentants en degrés vers les radians en conservant tant
les arguments des canaux que les coefficients des fonctionnelles.
Cette précision permet de suivre la dépendance de l’action sans
inversion rétrospective de ses coordonnées.

La transition hexadique--octadique apporte le contenu propre du
fonctionnel de Barbero--Immirzi. Les espaces d'incidence fixent les
cardinaux ; le tour orienté fixe les multiplicités ; l'évaluation
des canaux détermine le scalaire. Son utilisation ultérieure dans
l'opérateur d'aire et dans l'action gravitationnelle rend visible la signification
de la structure produisant la valeur. Le paramètre remplit ainsi une
fonction démonstrative dans la relation entre incidence, orientation et
géométrie.

La construction CKM détermine le mélange au moyen d’un lecteur explicite. Le
lecteur rationnel transporte les coordonnées angulaires précédentes vers trois
angles de mélange et fixe leur phase ; son inverse reconstruit les coordonnées
de départ. La factorisation unitaire explicite les amplitudes et
l’invariant de Jarlskog conserve l’information de violation CP sous les
changements de phase admissibles. Sa comparaison quantitative s’effectue avec
la convention angulaire et la référence physique déclarées. La construction
leptonique se présente dans ses repères spectraux propres, en conservant la
différence entre les deux transports d’espèce.

Les résultats de mélange sont exprimés numériquement par
\begin{align*}
 (\theta_{12},\theta_{23},\theta_{13})
   &=(13.003313589509,\ 2.398056157332,\ 0.213977382244)^\circ,\\
 \delta_{\rm CP}&=65.676173123554^\circ,\qquad
 J_{\rm CKM}=3.1189723326632974\times10^{-5}.
\end{align*}
Le tableau~\ref{tab:ii-ckm-comparacion} compare leurs sinus, phase et
invariant au PDG 2025 dans la même paramétrisation : les cinq
écarts marginaux normalisés satisfont $|z|<0.029$.
Le résultat relie les coordonnées engendrées à des amplitudes unitaires
et à une mesure de phase qui reste invariante sous reparamétrisation.

La réalisation angulaire dans l’incidence de Witt établit un transport
isométrique explicite, normalisé polairement, avec un inverse sur son image.
L’égalité métrique du corollaire~\ref{cor:ii-witt-composed-reader}
précise sa fonction : la composition avec le lecteur angulaire conserve la
métrique induite, distingue des coordonnées angulaires différentes et permet
de reconstruire les trois coordonnées qui alimentent les règles sectorielles.
La comparaison entre le registre et l’incidence s’effectue ainsi au moyen
d’une application fidèle. La représentation déterminant conserve
l’orientation ; les changements de jauge transforment conjointement les données,
selon \ref{prop:ii-witt-gauge-covariance} ; la réalisation complexe conserve
la phase et l’ordre des facteurs unitaires de mélange. Sont différenciés
la sélection des opérations sectorielles, leur évaluation et le transport
de leur résultat, sans identifier le transport isométrique ultérieur à la sélection des
règles sectorielles qu’il évalue.

\Needspace{11\baselineskip}
\subsection{Constante de Boltzmann et énergie par décision tritique}

L’action et la période produisent
\(\epsilon_a=h_a/T_{108}=\hbar_a/t_P\) ; l’information du triplet
uniforme fixe l’énergie par nat \(\epsilon_a/\log3\).
Une section thermique positive \(\Theta_{{\rm clk},a}\) étant fixée,
il existe un unique transducteur linéaire positif
\(k_B:\mathcal L_\Theta\to\mathcal L_E\) qui la transforme en cette
énergie. Son expression dans des bases positives conserve la loi
\(b'=(s_\Theta/s_E)b\). Sont ainsi déterminées tant la fonction
dimensionnelle de Boltzmann que les conditions qui individualisent sa
coordonnée : section thermique, bases et, lorsqu’une sélection
interne est revendiquée, le lecteur qui les fixe.

Le même transducteur agit sur les deux orientations lorsque les
températures du cycle ont le quotient des sections d’action,
selon \eqref{eq:ii-kb-aut-orientaciones}. Il assigne à l’information
sa grandeur entropique au moyen de \(S=k_Bs\), sans identifier une distribution
de continuations à une bipartition quantique ni une mise à jour
réversible à un effacement physique. L’énergie, l’information et
l’entropie conservent ainsi leurs domaines dans leurs relations avec l’aire et
avec la cellule gravitationnelle.

La construction de capacité de \ref{sec:ii-ter27-antecedentes} conserve le calendrier
de capacité avec ses vacances, les deux graduations du porteur et la
réalisation énergétique positive de sa mémoire. Le spectre décimal,
l’origine de capacité et les multiplicités déterminent l’intensité
\(b_*=1.380649\times10^{-23}\) dans l’assignation des degrés explicitée.
Les courants fixent le décalage dans sa réalisation connexe ; le transport
centré conserve le lecteur avec sa référence.

Le théorème~\ref{thm:ii-ter27-transducer} ajoute une conséquence
thermométrique précise. La référence fixe et l’état d’essai produisent
un accroissement entropique \(b_*s\) ; sa lecture conjointe avec l’énergie
conditionnée détermine \(\Theta_P=1/(b_*\beta)\) et un unique
transducteur positif de coefficient \(b_*\). Le
corollaire~\ref{cor:ii-ter27-variational} démontre en outre le minimum
unique d’énergie libre. La référence fixe et la paire de fonctionnelles
constituent les conditions de cette réalisation. Dans le porteur gradué,
le hamiltonien chronologique de capacité donne exactement
\(\rho=q^N/Z_N(q)\) ; son lien avec l’aire et l’horizon conserve les
échelles distinctes et l’identité finie d’entropie relative.

La section~\ref{sec:ii-aditiva-instrumento} incorpore la caractérisation
additive de ce poids et la réalisation conservative de ses marques.
La réunion d’alternatives disjointes, leur mesure reçue et l’unité de
référence fixent la contribution $b_*s$ ; le conditionnement retrouve
l’énergie moyenne du même essai. L’instrument conserve le complément
spectral, l’énergie du système et le transport tritique avec quotient et retenue.
Une référence partagée maintient un unique facteur $b_*$ lors de la composition de
systèmes et conserve leurs corrélations. Ces preuves explicitent les
opérations de la réalisation thermique, outre l’évaluation de son coefficient.

\subsection{Gravitation et conservation de l'échelle aréale}

La construction gravitationnelle conserve également une information tensorielle
que la section scalaire de couplage n’épuise pas. La descente d’incidence
de \ref{sec:ii-pentadica}, sous ses conditions explicites, produit le
porteur à cinq composantes. L’entrelaceur \(\Gamma_5\) réalise ce
porteur comme tenseurs symétriques sans trace ; sa base distingue les poids
\(\pm2,\pm1,0\). La réduction transversale en conserve deux au moyen
d’un projecteur de rang deux. Cette séquence transmet simultanément le
couplage commun \(G_a\), l’orientation de ses cinq composantes et
l’observation réduite, conformément à
\eqref{ii:pent:eq:gravity-12-5-5-2}. L’interprétation physique conserve la
dynamique et les contraintes spécifiées dans ce développement.

L’inscription, l’archive complète et les compositions longitudinale
et métrique déterminent d’abord le lecteur radial. L’énergie et le rayon
agissent sur le même caractère positif ; leur annulation démontre que
le coefficient de couplage est commun à tous les modes du domaine
déclaré. Les règles R4 et R5 figurent comme compositions incorporées,
et R8 maintient l’identification physique du rayon réduit. L’horloge
\(t_0=\pi_{\HMT}L/(54c_{\rm int})\) exprime ensuite la
même longueur dans le calendrier circulaire.

Le transport unitaire conserve le coefficient et les registres.
L’élimination conjointe de la frontière et de la mémoire conserve la proportion
entre rayon et énergie, le produit radial et le résidu récupérable ;
dans la résolvante complète, on transporte également le paramètre spectral,
selon \eqref{g26:eq:resolvente-proporcional-es}. Cette conservation
couvre la mémoire statique et la dynamique dans leurs domaines.
La réciprocité des sections d’action conserve le
produit de la gravitation par l’action et, avec lui, l’échelle d’aire. La
réalisation thermique relie l’énergie de la durée élémentaire à
l’information tritique. Dans la cellule d’horizon spécifiée, l’énergie,
la température, l’entropie et la demi-action satisfont les identités
démontrées. L’article situe ainsi la relation entre information et
gravité dans une composition mathématique explicite.

\subsection{Information de frontière et entropie d'intrication}

La lecture conjointe de l'aire et de l'information sectorielle fournit une
expression concrète de la généalogie de frontière. L'aire agrège les
occupations, tandis que l'observable pondérée conserve leur distribution
entre les deux secteurs. Leur récupération est démontrée par
l'inversion de l'application d'incidence. L'entropie d'intrication ajoute
l'état réduit et la bipartition qui lui donnent sens. Le
transport de ces données permet de distinguer conservation spectrale,
conservation de l'orientation et conservation de l'information accessible
dans une réduction.

La formulation sectorielle obtient en outre les moyennes de l'aire et de
l'information modulaire à partir de la fonction de partition finie. Ses
dérivées s'expriment par des variances et des covariances d'occupation.
L'inversion unitaire des occupations conserve l'entropie et
transforme l'aire en son complément. Se trouve ainsi précisé un phénomène
central de la lecture généalogique : une quantité conservée peut
coexister avec une transformation structurelle qu'une autre observable
détecte. Le transducteur de Boltzmann conserve sa fonction dimensionnelle
dans cette relation entre information, énergie et entropie.

Le transport des modes collectifs explicite comment l'incidence
hexadique--octadique atteint les étiquettes aréales. La première loi
modulaire relie, à état de référence fixé, les variations de
l'entropie et de l'occupation pondérée. Pour des changements finis, l'identité
inclut l'entropie relative entre les états ; conserver ce terme
précise le contenu informationnel distinguant une variation
infinitésimale d'une transformation finie. La récupération conjointe
des occupations par l'aire et l'incidence situe cette relation
directement sur les deux observables de frontière.

\subsection{Réciprocité et continuation exceptionnelle}

L'ellipse d'action et la dualité T prolongent l'argument par
des invariants de transport. L'échange des demi-axes conserve le
déterminant de la forme ; l'orientation symplectique distingue des
réalisations de même énergie ; l'échange moment--enroulement
conserve l'opérateur spectral sous l'identification correspondante
des domaines. La réalisation exceptionnelle maintient l'incidence
jusqu'à Moonshine et permet de formuler la réciprocité modulaire dans son
espace propre.

La réalisation de covariance de la
section~\ref{sec:ii-area-covarianza} précise le contenu opératoire de
cette réciprocité. Dans la normalisation de deux écarts-types,
l'aire \(h\) fixe
\(\det V_\eta=\hbar^2/4\), tandis que l'anisotropie angulaire détermine
\(\Delta P/\Delta Q=e^{-\eta}\). L'état gaussien construit et ses
opérateurs canoniques réalisent exactement ces deux relations ; la
rotation symplectique échange leurs dispersions tout en conservant le
commutateur. Se trouvent ainsi reliés le déterminant d'action,
la géométrie de l'ellipse et la saturation de l'incertitude dans une
représentation explicite. La distinction entre l'énergie du contour,
\(\hbar\omega\), et l'énergie moyenne de l'état,
\(\hbar\omega/2\), préserve la normalisation des deux lectures.

Le prolongement à la théorie M nécessite de spécifier la compactification,
les surfaces d’univers et les représentations physiques correspondantes.
La relation particule--antiparticule au moyen de l’orientation, de l’avance et du
retard constitue une proposition d’extension, distincte des
identités de covariance et de réciprocité démontrées ici.

La section~\ref{delta22:ii:transporte} démontre la conservation exacte de l'action quadratique sous transport entre ellipses, dérive son terme de connexion hamiltonienne et établit la condition de positivité du générateur. La version conformément symplectique conserve l'action normalisée et explicite la forme transportée.

\subsection{Évaluation du registre et conservation de l'observation}

L’évaluation d’incidence de
\eqref{eq:registro-composicion-incidencial} rend explicite l’opération
qui transforme un registre de visites et des traces orientées en vingt-cinq
coordonnées. Chaque triplet de dénombrements est représenté par ses résidus
décimaux, tandis que ses quotients et provenances demeurent dans le registre.
La proposition~\ref{prop:registro-36-residuos} détermine quand deux
familles de dénombrements ont la même image. Son contenu permet de
préciser quelle information utilise la valeur obtenue et quelle information
supplémentaire conserve l’histoire qui l’a produite. La famille d’incidences
et les traces qui alimentent l’évaluateur conservent leurs conditions de
construction ; l’inversibilité ultérieure du registre a une fonction
complémentaire, la récupération exacte de cette représentation.

La conservation de l'observation acquiert une formulation dynamique dans
\ref{subsec:ii-memoria-resolvente}. L'identité finie retient l'état
terminal de l'histoire ; son prolongement contrôle la contribution des
incréments ultérieurs par une majoration du reste. Lors de l'élimination
d'un secteur auxiliaire, les identités
\eqref{mem:schur-fuente}--\eqref{mem:schur-lector} transportent aussi
la source et l'application de lecture. Le contenu d'une trajectoire
peut ainsi s'exprimer en moins de composantes tout en conservant les effets de
ses excursions et de ses retours. La transitivité de
\eqref{mem:schur-transitivo} assure la cohérence de cette opération
lorsque l'élimination s'effectue en plusieurs étapes. Il s'agit d'une
conséquence précise du traitement de la mémoire comme partie opératoire de l'état.

% BEGIN THERMAL_R03_SYNC_13B
% Formalización reunida de resultados ya expuestos en el cuerpo.
\subsection{Définitions structurelles et résultats quantitatifs}
\label{sec:ii-conclusiones-cuantitativas}

La définition par la structure conserve l’objet antérieur à la lecture
scalaire. Pour les constantes d’action, cet objet est une section orientée
avec une norme multiplicative sur le quotient local et un caractère régional.
Pour Barbero--Immirzi, c’est
l’évaluation d’une incidence orientée. Pour le mélange, c’est un transport
entre coordonnées et amplitudes. Pour Boltzmann, c’est une application entre
droites de grandeurs. Pour la gravitation, c’est une section réciproque de
l’action dont la grandeur provient de la longueur inscrite et de la lecture
radiale spécifiées dans R1--R8. Ces définitions
déterminent des opérations différentes dans la même généalogie.

Le tableau suivant fixe, pour chaque grandeur, l’objet qui la définit,
l’opération qui détermine sa lecture et la fonction physique développée
dans le corps du texte. Les preuves de reconstruction et de transport établissent
comment cette information est conservée lors de la réutilisation du résultat. En
particulier, \ref{prop:registro-36-residuos} caractérise la représentation
du registre, \ref{mem:prop-schur} conserve l’observation sous réduction
et \ref{cor:ii-witt-composed-reader} conserve les coordonnées angulaires
dans la réalisation d’incidence. Ces trois opérations ont des domaines
propres et des fonctions démonstratives distinctes au sein de la même séquence.

\begingroup
\setlength{\LTpre}{0pt}
\setlength{\LTpost}{2pt}
\begin{longtable}{@{}p{.13\linewidth}p{.415\linewidth}p{.395\linewidth}@{}}
\caption{Objet structurel, lecture produite et sens physique des
résultats de l’article.}\label{tab:ii-conclusiones-estructura}\\
\toprule
Grandeur & Définition structurelle et résultat & Lecture physique et développement interne\\
\midrule
\endfirsthead
\toprule
Grandeur & Définition structurelle et résultat & Lecture physique et développement interne\\
\midrule
\endhead
\bottomrule
\endfoot
Planck & Section de la droite d’action :
$\hbar_{\rm pre}=H_5\,10^{-34}\mathcal U_S$,
$\hbar_{\rm ret}=\eta_{\rm ret}\,10^{-34}\mathcal U_S$.
Le quotient local a seize feuillets et
$\operatorname{ord}_{10}(\varphi\alpha^{16})=-34$.
& La section relie la fréquence et l’énergie ; le tour complet détermine
$h_a=2\pi\hbar_a$. Développement dans \ref{sec:action} ; coordonnées et comparaison
dans \ref{tab:ii-planck-comparacion}.\\[5pt]
Barbero--Immirzi & Évaluation orientée des canaux et de l’incidence :
$\gamma_{\HMT}=\pi AC^*/180+12\Delta S_{90}-\Delta S_{120}$.
L’évaluation imprimée est $0.2740076726944592747\ldots$.
& Le coefficient intervient dans la résolution élémentaire d’aire et dans
la réalisation Cartan--Holst. Les définitions de
\ref{sec:barbero} et \ref{sec:area} conservent ces deux usages ;
\eqref{eq:holst} spécifie le second.\\[5pt]
Mélange CKM & Image du vecteur angulaire par les trois applications
sectorielles ; phase $\delta_{\deg}=9A_{\deg}$.
Les angles et les amplitudes sont évalués après la construction du lecteur.
& Le quartet de Jarlskog exprime l’information de phase invariante
sous les changements diagonaux de base. Les règles, leur évaluation et la
réalisation unitaire sont développées dans
\ref{sec:ii-propagacion-angular}.\\[5pt]
Boltzmann & Transduction positive
$k_B:\mathcal L_\Theta\to\mathcal L_E$ avec
$k_B\Theta_{{\rm clk},a}\log3=h_a/T_{108}$.
L’état et la mesure des continuations déterminent le contenu
informationnel sur lequel elle agit.
& $S=k_Bs$ réalise dimensionnellement l’entropie de l’état réduit.
Les bases d’énergie et de température et leurs transformations sont conservées
dans \ref{sec:ii-boltzmann} ; la caractérisation additive et l’instrument
sont démontrés dans \ref{sec:ii-aditiva-instrumento} ; la comparaison figure dans
\ref{tab:ii-kb-g-cartas}.\\[5pt]
Réalisation thermique marquée & Référence fixe $\eta$, avec les multiplicités
$(d_0,d_1,d_2)=(1,380,649)$ aux degrés $23,26,29$ et la trace
$b_*=1.380649\times10^{-23}$.
Les lectures $\mathcal S_P=b_*s$ et $\mathcal E_P=\operatorname{Tr}(\rho H)$
déterminent $d\mathcal S_P/d\mathcal E_P=b_*\beta$ et
$\Theta_P=(b_*\beta)^{-1}$.
& Avec une référence fixe, des bases positives transportées et un hamiltonien non
scalaire, \ref{thm:ii-ter27-transducer} détermine le transducteur.
Le corollaire~\ref{cor:ii-ter27-variational} prouve le minimum unique.
La composition conserve $q^N/Z_N(q)$, l’origine énergétique et les
transports horloge--horizon de \eqref{eq:ii-ter27-intertwiner}.\\[5pt]
Gravitation & Section réciproque
$G_a=c_{\rm int}^{3}L^2/\hbar_a$,
avec $L=(5759/23040)\alpha_{\HMT}^{16}L_*$.
L’horloge $t_0=\pi_{\HMT}L/(54c_{\rm int})$ conserve l’expression
circulaire $G_a=\vartheta_{108}^{\circ}c_{\rm int}^{5}t_0^2/\hbar_a$,
avec $\vartheta_{108}^{\circ}=(54/\pi_{\HMT})^2$.
& Le produit $G_a\hbar_a$ conserve la normalisation aréale.
La construction longitudinale et radiale, ses règles et l’unicité du
coefficient sont démontrées dans \ref{g26:sec:rigidez-radial-es} ;
la réalisation circulaire est conservée dans \ref{sec:gravity} et la confluence thermique dans
\ref{sec:confluencia}.\\
\end{longtable}
\endgroup

La comparaison quantitative conserve le sens propre de chaque
référence. Pour Planck, le tableau~\ref{tab:ii-planck-comparacion} compare
des coordonnées dimensionnelles à une constante définissant le SI ; les
résidus relatifs des orientations antérieure et postérieure sont,
respectivement, d’environ $-1.82\times10^{-15}$ et
$-6.93\times10^{-10}$. Pour Barbero--Immirzi, le
tableau~\ref{tab:ii-barbero-comparacion} compare deux constructions :
l’évaluation interne et la racine de l’équation de Ghosh--Mitra, dont la
valeur approchée est $0.2740668582328300$. La différence relative est
d’environ $-2.16\times10^{-4}$ ; l’objet comparé est un
coefficient théorique.

Pour CKM, le tableau~\ref{tab:ii-ckm-comparacion} réunit les sinus des
trois angles de mélange, la phase en radians et l’invariant de Jarlskog.
Les angles et la matrice des modules sont explicités dans le développement
adjacent.
Les écarts marginaux normalisés qui y sont définis satisfont
$|z|<0.029$ par rapport aux lignes PDG utilisées. Chaque résidu utilise
l’incertitude marginale indiquée. Leur lecture est individuelle ;
une évaluation conjointe nécessite en outre la matrice de covariance
de ces observables.
Pour Boltzmann et la gravitation, le tableau~\ref{tab:ii-kb-g-cartas}
expose séparément le transducteur, la composition longitudinale et radiale, les coordinatisations dimensionnelles
et la référence externe. Cette organisation permet
d’identifier exactement la quantité que produit chaque construction et
celle qu’évalue chaque comparaison.

Le tableau~\ref{tab:ii-conclusiones-contraste} réunit les résultats de
ces comparaisons avec les définitions structurelles précédentes.
\begin{table}[!htbp]
\centering\small
\setlength{\tabcolsep}{4pt}
\renewcommand{\arraystretch}{1.12}
\begin{tabularx}{\linewidth}{@{}>{\raggedright\arraybackslash}p{.16\linewidth}>{\raggedright\arraybackslash}X>{\raggedright\arraybackslash}X@{}}
\toprule
Grandeur & Résultat dans la réalisation spécifiée & Référence et comparaison\\
\midrule
Planck, avant le retour
& \(\hbar_{\rm pre}=1.0545718176461544696\ldots\)
  \(\times10^{-34}\,\mathrm{J\,s}\).
& SI : \(\hbar_{\rm SI}=h_{\rm SI}/(2\pi)\), exact.
  Sa coordonnée est \(1.0545718176461563912\ldots\)
  \(\times10^{-34}\,\mathrm{J\,s}\).
  Résidu relatif \(-1.822221\times10^{-15}\).\\[4pt]
Planck, après le retour
& \(\hbar_{\rm ret}=1.0545718169150390611\ldots\)
  \(\times10^{-34}\,\mathrm{J\,s}\).
& Même référence SI.
  Résidu relatif \(-6.932836\times10^{-10}\).
  Pour \(h_a=2\pi\hbar_a\), les résidus sont identiques.\\[4pt]
Barbero--Immirzi
& \(\gamma_{\HMT}=0.2740076726944592747\ldots\).
& Ghosh--Mitra : environ \(0.2740668582328300\).
  Différence relative \(-2.159529202\times10^{-4}\)
  entre constructions théoriques.\\[4pt]
Mélange CKM
& \(\delta_{\rm CP}=65.676173123554^\circ\),
  \(J=3.1189723326632974\times10^{-5}\).
  Les trois angles de \eqref{eq:ii-ckm-decimales} sont
  \(\theta_{12}=13.003313589509^\circ\),
  \(\theta_{23}=2.398056157332^\circ\) et
  \(\theta_{13}=0.213977382244^\circ\).
& PDG 2025 : les cinq comparaisons de
  \(s_{12},s_{23},s_{13},\delta,J\) satisfont \(|z_x|<0.029\)
  avec les incertitudes marginales du
  tableau~\ref{tab:ii-ckm-comparacion}.\\[4pt]
Boltzmann
& \(B_*=1.380649\times10^{-23}\,u_E/u_\Theta\),
  avec les bases et la règle thermométrique de
  \eqref{eq:ii-b-evaluacion-comparada}.
& SI : \(1.380649\times10^{-23}\,\mathrm{J/K}\), exact.
  Coïncidence du nombre sous le transport dimensionnel déclaré.\\[4pt]
Longueur \(L\)
& \(L=1.6162549826251263301\ldots\)
  \(\times10^{-35}\,\mathrm m\), pour \(L_*=1\,\mathrm m\).
& Composition longitudinale et radiale :
  \(G_{\rm ret}=c_{\rm int}^3L^2/\hbar_{\rm ret}\),
  avec les mêmes bases de vitesse et d’action.\\[4pt]
Gravitation
& \(G_{\rm ret}=6.6742996594140227\ldots\)
  \(\times10^{-11}\,\mathrm{m^3\,kg^{-1}\,s^{-2}}\).
& CODATA 2022 : \(6.67430(15)\times10^{-11}\) dans les mêmes unités.
  Différence d’environ \(-0.00227057\) incertitudes-types.\\
\bottomrule
\end{tabularx}
\caption{Synthèse de la comparaison quantitative. Les évaluations, les bases
et les références sont celles des tableaux~\ref{tab:ii-planck-comparacion}--\ref{tab:ii-kb-g-cartas}.
Les résidus relatifs des constantes d’action, la comparaison
théorique de Barbero--Immirzi et les résidus normalisés de CKM et de la gravitation
conservent leurs sens respectifs.}
\label{tab:ii-conclusiones-contraste}
\par\medskip
\centering\small
\setlength{\tabcolsep}{4pt}
\renewcommand{\arraystretch}{1.18}
\begin{tabular}{@{}lrrr@{}}
\toprule
Observable & Évaluation HMT & PDG 2025 & \(z_x\)\\
\midrule
\(s_{12}\) & \(0.225007404757\) & \(0.22501\pm0.00068\) & \(-0.003817\)\\
\(s_{23}\) & \(0.041841757010\) & \(0.04183^{+0.00079}_{-0.00069}\) & \(0.014882\)\\
\(s_{13}\) & \(0.003734601164\) & \(0.003732^{+0.000090}_{-0.000085}\) & \(0.028902\)\\
\(\delta\) (rad) & \(1.146265461116\) & \(1.147\pm0.026\) & \(-0.028251\)\\
\(10^5J\) & \(3.118972333\) & \(3.12^{+0.13}_{-0.12}\) & \(-0.008564\)\\
\bottomrule
\end{tabular}
\caption{Valeurs CKM et comparaison marginale réunies dans les conclusions.
Les trois angles du tableau~\ref{tab:ii-conclusiones-contraste} sont comparés
au moyen de \(s_{ij}=\sin\theta_{ij}\), avec la phase en radians et
l’invariant \(J\). On conserve les chiffres, incertitudes et résidus du
tableau~\ref{tab:ii-ckm-comparacion}, correspondant à l’ajustement à trois
générations PDG 2025. La dernière ligne exprime conjointement la valeur et
l’incertitude en unités de \(10^{-5}\). Les résidus sont des comparaisons
marginales d’observables corrélées.}
\label{tab:ii-conclusiones-ckm}
\end{table}


\Needspace{10\baselineskip}
\subsection{Conservation de l’échelle aréale et récupération de l’information de frontière}

La relation entre information et gravitation se concentre dans deux applications
complémentaires. La première conserve l’échelle aréale au moyen de
$G_a\hbar_a$ et organise ses occupations au moyen de
$\gamma_{\HMT}a_0$. La seconde conserve la résolution sectorielle au moyen de
l’observation conjointe de l’aire et de l’incidence pondérée. L’inversion
\eqref{eq:recover} retrouve les occupations à partir de l’aire et de
l’observable pondérée ; les équations
\eqref{eq:ii-respuesta-entropica} et \eqref{eq:ii-respuesta-areal}
relient leurs variations aux fluctuations de l’état.

De là provient la lecture physique centrale de l’article : une même
configuration admet une observation géométrique et une observation
informationnelle reliées par une structure récupérable. Le changement
d’orientation peut conserver l’entropie de l’état réduit et
transformer l’aire, comme l’établissent
\eqref{eq:ii-inversion-entropia} et \eqref{eq:ii-area-complementaria}.
La conservation de l’information se rapporte à l’état et à son transport ;
l’entropie d’intrication se rapporte en outre à la bipartition et à
la réduction. L’identité d’horizon de
\ref{sec:confluencia} utilise ces fonctions dans la réalisation
géométrique et thermique spécifiée.

Les identités de mémoire permettent également de préciser la portée d’une
réduction de variables. L’égalité des observations de
\eqref{mem:schur-lector} conserve les contributions du secteur éliminé
au moyen de l’opérateur effectif, de la source et du lecteur transformés.
L’entropie de l’état réduit, en revanche, s’évalue avec la bipartition
et la mesure spécifiées dans \ref{sec:ii-ecuacion-informacion}.
Les deux constructions expliquent des fonctions différentes de la mémoire :
la première conserve une lecture du transport ; la seconde quantifie
l’information accessible au sous-système. Leur articulation maintient
explicites la structure transportée et la réduction qui définit
l’observable physique.

Le progrès explicatif s’exprime donc au moyen d’un ordre précis :
structure engendrée, application de lecture, valeur obtenue et fonction
physique. L’égalité de deux valeurs s’étudie avec l’égalité
ou la différence de leurs structures ; les applications de récupération déterminent
laquelle de ces différences conserve un sens observable.
% END THERMAL_R03_SYNC_13B


Les comparaisons quantitatives maintiennent la distinction entre une
constante définissant des unités, un coefficient théorique et des observables
de mélange assorties d’incertitudes. Elles examinent les coordonnées produites dans
les coordinatisations déclarées, tandis que les identités structurelles déterminent
ce qui est conservé lors de leur transport : le quotient de retour de l’action,
l’orientation de la fonctionnelle de Barbero--Immirzi, la métrique et la phase
du mélange, la transduction thermique et le produit gravitationnel.

La contribution de l’article est cette détermination articulée des
constantes et de leurs relations physiques. La composition APP--TRIT--TPK
construit les régions numériques et la mémoire dodécaphasique avant leurs
évaluations ; les lecteurs ultérieurs produisent les sections d’action,
les coordonnées angulaires et les observables de frontière. Dans la réalisation
d’horizon spécifiée, \(T_HS_H=E_\Lambda\) et la cellule de rayon
\(\ell_P\) réunit l’énergie de décision et la demi-action. L’action fixe
l’énergie par durée ; Barbero--Immirzi résout l’incidence aréale ;
Boltzmann relie l’énergie et l’information ; et la gravitation conserve
l’unité géométrique \(G_a\hbar_a/c_{\rm int}^3\).
Ces relations expriment le sens des constantes au sein
d’une structure commune et conservent les conditions précises sous lesquelles
chacune est déterminée.
% END THERMAL_R03_SYNC_13
